% SPDX-License-Identifier: CC-BY-4.0
\section{Introduction}

Lunar PNT has stopped being a slide and started being a procurement. ESA's
Moonlight initiative is now building a lunar communications and navigation
service (LCNS) from a small constellation, one communications and four
navigation satellites in elliptical lunar frozen orbits; the
EUR~123~million Phase-1 prime contract was signed in October 2024 and the space
segment is in industrial build \citep{esa_moonlight,telespazio2024moonlight}. At
its CM25 Ministerial in November 2025 ESA also adopted NovaMoon, a
geodetic-timing-differential reference station that puts multi-technique geodesy
(VLBI and ranging), joint orbit-and-clock determination, reference-frame
realisation, timekeeping, and differential corrections on the lunar surface
\citep{molli2026novamoon}. NASA, ESA and JAXA have agreed a common
interoperability framework, the LunaNet Interoperability Specification (LNIS~v5,
January 2025), and all three broadcast its Augmented Forward Signal
\citep{nasa2025lnis,nasa2025afs}. A Coordinated Lunar Time exists on paper too:
the IAU's 2024 General Assembly defined a lunar coordinate time scale (TCL)
\citep{iau2024lcrs}, a NASA-led realisation is due by the end of 2026
\citep{ostp2024lunartime}, and the relativistic rate (roughly
56--57~$\mu$s/day faster than terrestrial time) is set by a published framework
\citep{ashby2024lunartime}. And in March 2025 LuGRE took the first GNSS fix on
the lunar surface \citep{lugre2026}. Underneath all of it sits one problem:
working out the orbits and clocks of the lunar infrastructure and tying them to a
terrestrial or inertial frame.

The first thing you ask about that problem is observability. Given a set of
measurement types, which states (positions, clocks, frame parameters) come out
well determined, and which do not? The answer is an architecture decision: does
a lunar nav service need Earth-baseline geodetic ties, how many, and how good do
they have to be? It is not even an ESA-specific decision. Because ESA's LCNS,
NASA's LCRNS and JAXA's LNSS all share the Augmented Forward Signal under
LNIS~v5, the same observability question applies to the whole stack
\citep{nasa2025lnis}. Yet it is usually answered in private, one-off studies
whose assumptions never leave the building. We are not aware of an open, citable
tool that regenerates a lunar PNT network's observability from a committed
scenario and labels each figure of merit as externally validated or modelled,
the discipline the kshana program has applied to terrestrial and deep-space PNT
\citep{baweja2026kshana} and that we extend here to the lunar case
\citep{iiyama2023lupnt,fernandezprades2018reproducibility}.

The contribution of this paper is an observability result and a design law that
follows from it. We prove that the absolute position of a lunar surface station
in joint orbit-and-clock determination can be fixed only through a tie to the
Earth/inertial frame, and that two such ties exist: an indirect one in which
Earth$\to$satellite ranging anchors the constellation and lunar-local ranging
propagates that anchor to the station, and a direct one in which Earth-baseline
VLBI ties the station beacon itself. The indirect tie reaches the station only
when the constellation geometry is rich enough; for a sparse or early
constellation the station's absolute position falls in the null space of the
Fisher information and has no finite error bound at all, until a direct VLBI tie
restores it. This yields a clean conditional design rule (Theorem~\ref{thm:rigid}
and Table~\ref{tab:restore}): an Earth-baseline VLBI tie, of the kind NovaMoon's
transmitter would provide, is \emph{necessary} when the constellation cannot
carry the absolute datum on its own and merely \emph{sharpens} the bound when it
can. The direct tie has its own threshold, derived from the geometry: because
Earth stations subtend a tiny angle from the Moon, a single VLBI epoch informs at
most two of the station's three absolute axes, so the absolute position becomes
observable at three non-collinear Earth stations and not before
(\eqref{eq:defect}, Figure~\ref{fig:n_earth}). The familiar large
improvement-factor headline---a $91\times$ ensemble-median reduction in absolute
station error---survives as one honest illustration of the sparse regime, where
the indirect tie has nothing to give and VLBI is the difference between an
unobservable datum and a metre-class one; it is the consequence of the
observability result, not its substance. Everything ships as an open,
reproducible toolkit, seven capability areas in the \code{kshana} engine, with
each figure of merit labelled validated or modelled.
